% GENERATED FILE. Edit canonical lessons, then run npm run generate:books.
% canonical-source-hash: fnv1a64:2e6a6e92f43d3b3d
% canonical-lessons: PT-C437-avenida, PT-C437-travessa, PT-C437-beco, PT-C437-viela, PT-C437-bifurcacao

\chapter{Ruas e becos}
\label{ch:pt-ruas-e-becos}
\hlchaptermodality{437}

\begin{chapteropening}
\textbf{By the end of this chapter:} \emph{I can describe streets: the avenue, the side street, the alley, the narrow lane and a fork in the road.}

Complete the last lesson of chapter 437: I can describe streets: the avenue, the side street, the alley, the narrow lane and a fork in the road.
\end{chapteropening}

\section[a avenida]{a avenida — the avenue}

\label{lesson:PT-C437-avenida}



\begin{quote}
Before the new one: say the Portuguese for the road junction, then the Portuguese for to carry on.
\end{quote}

\subsection*{You'll want to know: a avenida}
\textbf{a avenida} — \textquotedblleft{}the avenue\textquotedblright{}.

\begin{quote}
A avenida principal tem muitas lojas.
\end{quote}

\begin{quote}
\emph{The main avenue has a lot of shops.}
\end{quote}

Shortened to \emph{Av.} in addresses.

\begin{cousinweb}
\emph{avenida} was borrowed from Spanish \emph{avenida}, which came from French \emph{avenue} 'way of approach' (EtymDB), on Latin \emph{advenīre} 'to arrive'.
\end{cousinweb}

\begin{grammarlens}[title={What you've built}]
One more word for telling, explaining and describing at B1.
\end{grammarlens}

\subsection*{Your turn}
\begin{itemize}
\raggedright
  \item \emph{Say it:} \emph{a avenida}
  \item \emph{Say it:} \emph{a avenida}, once more
  \item \emph{Say it:} a sentence of your own with \emph{a avenida}
  \item \emph{Recall:} say the Portuguese for the road junction, then the Portuguese for to carry on, then say \emph{a avenida} again
\end{itemize}

\begin{tcolorbox}[breakable,colback=teal!4,colframe=teal!35!black,title={Before you move on}]
What does a avenida mean? (\textquotedblleft{}The avenue\textquotedblright{}.) Say it once more.
\end{tcolorbox}
\section[a travessa]{a travessa — the side street; the serving dish}

\label{lesson:PT-C437-travessa}



\begin{quote}
Before the new one: say the Portuguese for to carry on, then the Portuguese for the avenue.
\end{quote}

\subsection*{You'll want to know: a travessa}
\textbf{a travessa} — \textquotedblleft{}the side street; the serving dish\textquotedblright{}.

\begin{quote}
A padaria fica numa travessa da rua principal.
\end{quote}

\begin{quote}
\emph{The bakery is in a side street off the main road.}
\end{quote}

For serving food, \emph{uma travessa} is a big dish.

\begin{cousinweb}
\emph{travessa} goes back to Latin \emph{trānsversa} 'lying across', from \emph{trānsvertere} 'to turn across' (Lewis \& Short): a street that crosses another.
\end{cousinweb}

\begin{grammarlens}[title={What you've built}]
One more word for telling, explaining and describing at B1.
\end{grammarlens}

\subsection*{Your turn}
\begin{itemize}
\raggedright
  \item \emph{Say it:} \emph{a travessa}
  \item \emph{Say it:} \emph{a travessa}, once more
  \item \emph{Say it:} a sentence of your own with \emph{a travessa}
  \item \emph{Recall:} say the Portuguese for to carry on, then the Portuguese for the avenue, then say \emph{a travessa} again
\end{itemize}

\begin{tcolorbox}[breakable,colback=teal!4,colframe=teal!35!black,title={Before you move on}]
What does a travessa mean? (\textquotedblleft{}The side street; the serving dish\textquotedblright{}.) Say it once more.
\end{tcolorbox}
\section[o beco]{o beco — the alley, the dead end}

\label{lesson:PT-C437-beco}



\begin{quote}
Before the new one: say the Portuguese for the avenue, then the Portuguese for the side street.
\end{quote}

\subsection*{You'll want to know: o beco}
\textbf{o beco} — \textquotedblleft{}the alley, the dead end\textquotedblright{}.

\begin{quote}
O carro entrou num beco sem saída.
\end{quote}

\begin{quote}
\emph{The car went into a dead end.}
\end{quote}

\emph{Um beco sem saída} — a dead end; also a hopeless situation.

\begin{cousinweb}
The origin of \emph{beco} is not settled, and EtymDB has no entry for it.
\end{cousinweb}

\begin{grammarlens}[title={What you've built}]
One more word for telling, explaining and describing at B1.
\end{grammarlens}

\subsection*{Your turn}
\begin{itemize}
\raggedright
  \item \emph{Say it:} \emph{o beco}
  \item \emph{Say it:} \emph{o beco}, once more
  \item \emph{Say it:} a sentence of your own with \emph{o beco}
  \item \emph{Recall:} say the Portuguese for the avenue, then the Portuguese for the side street, then say \emph{o beco} again
\end{itemize}

\begin{tcolorbox}[breakable,colback=teal!4,colframe=teal!35!black,title={Before you move on}]
What does o beco mean? (\textquotedblleft{}The alley, the dead end\textquotedblright{}.) Say it once more.
\end{tcolorbox}
\section[a viela]{a viela — the narrow lane, the alleyway}

\label{lesson:PT-C437-viela}



\begin{quote}
Before the new one: say the Portuguese for the side street, then the Portuguese for the alley.
\end{quote}

\subsection*{You'll want to know: a viela}
\textbf{a viela} — \textquotedblleft{}the narrow lane, the alleyway\textquotedblright{}.

\begin{quote}
No bairro antigo, as vielas são estreitas e escuras.
\end{quote}

\begin{quote}
\emph{In the old quarter, the alleyways are narrow and dark.}
\end{quote}

Typical of the old quarters of hill towns.

\begin{cousinweb}
\emph{viela} is usually explained as a small form of \emph{via} 'way' with \emph{-ela}, from Latin \emph{via} (Lewis \& Short); EtymDB has no Portuguese entry.
\end{cousinweb}

\begin{grammarlens}[title={What you've built}]
One more word for telling, explaining and describing at B1.
\end{grammarlens}

\subsection*{Your turn}
\begin{itemize}
\raggedright
  \item \emph{Say it:} \emph{a viela}
  \item \emph{Say it:} \emph{a viela}, once more
  \item \emph{Say it:} a sentence of your own with \emph{a viela}
  \item \emph{Recall:} say the Portuguese for the side street, then the Portuguese for the alley, then say \emph{a viela} again
\end{itemize}

\begin{tcolorbox}[breakable,colback=teal!4,colframe=teal!35!black,title={Before you move on}]
What does a viela mean? (\textquotedblleft{}The narrow lane, the alleyway\textquotedblright{}.) Say it once more.
\end{tcolorbox}
\section[a bifurcação]{a bifurcação — the fork (in the road), the split}

\label{lesson:PT-C437-bifurcacao}



\begin{quote}
Before the new one: say the Portuguese for the alley, then the Portuguese for the narrow lane.
\end{quote}

\subsection*{You'll want to know: a bifurcação}
\textbf{a bifurcação} — \textquotedblleft{}the fork (in the road), the split\textquotedblright{}.

\begin{quote}
Na bifurcação, a estrada da direita vai para a praia.
\end{quote}

\begin{quote}
\emph{At the fork, the right-hand road goes to the beach.}
\end{quote}

The verb is \emph{bifurcar-se}.

\begin{cousinweb}
\emph{bifurcação} is a learned formation on Latin \emph{bifurcus} 'two-pronged', from \emph{bi-} 'two' and \emph{furca} 'fork' (Lewis \& Short).
\end{cousinweb}

\begin{grammarlens}[title={What you've built}]
One more word for telling, explaining and describing at B1.
\end{grammarlens}

\subsection*{Your turn}
\begin{itemize}
\raggedright
  \item \emph{Say it:} \emph{a bifurcação}
  \item \emph{Say it:} \emph{a bifurcação}, once more
  \item \emph{Say it:} a sentence of your own with \emph{a bifurcação}
  \item \emph{Recall:} say the Portuguese for the alley, then the Portuguese for the narrow lane, then say \emph{a bifurcação} again
\end{itemize}

\begin{tcolorbox}[breakable,colback=teal!4,colframe=teal!35!black,title={Before you move on}]
What does a bifurcação mean? (\textquotedblleft{}The fork (in the road), the split\textquotedblright{}.) Say it once more.
\end{tcolorbox}
